\scanpage{9}
\uline{Compléments} $\overline{\text{\uline{XII}}}$

6. \uline{Le tore central maximal}. \uline{Relations avec le plus grand Tore quotient d'un \unsure{gpe|groupe} algébrique}

6.1. \uline{Définition en général} \unsure{Compatible|Compatibilité} aux extensions de la base descente, \unsure{générisation|spécialisation}\ldots\ etc\ldots

6.2. \uline{Existence lorsqu'il existe un \unsure{centr.|cart.|autre} réductif}

6.3. \deleted{\uline{N.B.} \unsure{En particulier, les groupes algébriques|En particulier, des groupes algébriques}} \uline{Un groupe algébrique : un \deleted{\unsure{tore}} maximal central} [N.B. le groupe algébrique n'a pas à être affine\ldots].

\deleted{\uline{6.4.} Soit $G$ un groupe algébrique \inserted{lisse \deleted{\unsure{connexe}}\unsure{connexe} affine} sur un corps parfait, \deleted{alors} soit $R$ son tore central maximal. Alors il existe un sous-groupe fini $\unsure{Z|3}$ de $R$, tel que \unsure{!} $\deleted{\unsure{R}}G/Z$ soit \deleted{facteur direct du} isomorphe au produit de $\deleted{G}R/Z$ et d'un groupe $H$. De plus, $H$ est tel que}

\margin{(\unsure{Rosenlicht})}
\uline{Th. 6.4} $X,Y$ \uline{géom. intègres} sur un corps $k$, munis de pts $a,b$ à valeurs dans $k$. Soit $\varphi:X\times Y\unsure{\mathord{-}|\to}G_m$ tel que $\deleted{\unsure{f}}\varphi(x,y)=1$, alors $\varphi$ est de la forme
\[
\varphi(x,y)=\varphi_1(x)\varphi_2(y)
\]
où $\deleted{\varphi_1(\varphi_2)}$ $\unsure{v|u}$ $(w)$ est un morphisme pointé $X\mathrel{-}G_m$ $(Y\mathrel{-}G_m)$.
% archive sheet number: 8
